\slajdid{09_1-s009}
\begin{frame}[label=09_1-s009]{Logička kola sa MOS tranzistorima}
\gusttekst\begin{columns}[c,onlytextwidth]\column{.44\textwidth}\centering
\begin{tikzpicture}[x=.8cm,y=.7cm,font=\sitneoznake,>=Latex]
\draw(0,0)rectangle(2,3.6);\draw(.4,0)--(.4,3.6);\draw(.4,1.65)--(2,1.65);\draw(.4,1.95)--(2,1.95);\node[align=center]at(1.2,2.7){Pull up\\network};\node[align=center]at(1.2,.8){Pull down\\network};\draw(1.05,3.6)--(1.05,4);\draw(.8,4)--(1.3,4)node[above left]{$V_{DD}$};\draw(1.05,0)--(1.05,-.4)node[ground]{};
\draw[->](1.85,1.8)--(2.65,1.8)node[right]{Izlaz};\draw(-.75,1.98)--(-.3,1.98);\draw(-.75,1.62)--(-.3,1.62);\draw[->](-.8,1.8)--(0,1.8);\node[left]at(-.8,1.8){Ulazi};
\end{tikzpicture}

\column{.54\textwidth}
Ako je PUN mreža realizovana samo sa otpornikom nazivaćemo takva kola sa pasivnim opterećenjem na izlazu.\par\bigskip
Ako je PUN mreža realizovana preko aktivnih elemenata, tranzistora, takva kola ćemo nazivati kolima sa aktivnim opterećenjem na izlazu.
\end{columns}\medskip
1. Pull up network (PUN) deo kola koji na izlazu obezbeđuje logičku jedinicu\par
2. Pull down network (PDN) deo kola koji na izlazu obezbeđuje logičku nulu\par\bigskip
Cilj nam je bio da PUN i PDN budu komplementarne, odnosno da kada jedna radi druga je isključenja.\par\medskip
Videćemo kod \alert{trostatičkih kola} situaciju da obe mreže isključujemo.\par\medskip
Isto tako kola koja imaju samo PDN – \alert{kola sa otvorenim drejnom}.
\end{frame}
